\documentclass[a4paper,12pt]{article}
\usepackage[margin=2.5cm]{geometry}
\usepackage{graphicx}
\usepackage{amsmath}
\usepackage{amssymb}
\usepackage{siunitx}
\usepackage{hyperref}

\title{Meander Line Antenna Simulation}
\author{Christian Fenger \& Jeppe Iversen}
\date{\today}

\begin{document}

\maketitle

\section*{Calculations based on paper}
\begin{align*}
d &= 0.16\lambda_g \\
s &= 0.41\lambda_g \\
L &= 0.70\lambda_g \\
w &= 0.05\lambda_g
\end{align*}

$$
\lambda_g = \frac{\lambda}{\sqrt{\epsilon_{ff}}}
$$
$$
\epsilon_{ff} = \frac{\epsilon_r+1}{2} + \frac{\epsilon_r-1}{2} \left( 1 + 10\frac{h}{w} \right)^{-1/2}
$$

Assuming a standard 4-layer PCB:
\[
h = 1.6~\text{mm}, \quad \epsilon_r = 4.5
\]

First we calculate $\epsilon_{ff}$ to get $\lambda_g$:
$$\lambda_g = \frac{\lambda}{\sqrt{\epsilon_{ff}}},$$
$$\lambda = \frac{c}{f} = \frac{3\cdot 10^8}{1060~\text{MHz}} = 28.3~\text{cm},$$
$$\epsilon_r = 4.5, \quad h = 1.6~\text{mm},$$
$$ \therefore \epsilon_{ff} = 3.73$$

Therefore:
\[
\lambda_g = 0.1465~\text{m} = 14.65~\text{cm} \\
\]

The paper mentions a necessary dimension correction $\Delta L$:
\[
\Delta L = 0.5 w, \qquad 
w = \frac{c}{2f_r} \sqrt{\frac{2}{\epsilon_r + 1}}
\]

The meander line dimensions before corrections are:
\begin{align*}
d &= 0.16\lambda_g = 0.02344~\text{m} \\
s &= 0.42\lambda_g = 0.06153~\text{m} \\
L &= 0.70\lambda_g = 0.10255~\text{m} \\
w &= 0.05\lambda_g = 0.007325~\text{m}
\end{align*}

Now we calculate $\Delta L$:
\[
\Delta L = 0.5 \left( \frac{3\cdot 10^8~\text{m/s}}{2 \cdot 1.06~\text{GHz}} \right) 
\sqrt{\frac{2}{4.5+1}} = 0.04266~\text{m}
\]

The corrected values are:
\begin{align*}
S_p &= s+\Delta L = 0.10419~\text{m} \\
d_p &= d+\Delta L = 0.0661~\text{m} \\
L_p &= L+\Delta L = 0.14521~\text{m}
\end{align*}

Now we calculate the HPBWs and the directivity:
\[
k_0 = \frac{\pi}{\sqrt{\epsilon_{ff}}L_p} = 11.20208
\]

\begin{align*}
\theta_E &= 2.45~\text{rad} = 140^\circ \\
\theta_H &= 1.867~\text{rad} = 107^\circ
\end{align*}

Therefore:
\[
D = \frac{41253}{\theta_E \theta_H} = 2.754 \approx 4.4~\text{dBi}
\]

\section*{Simulation Results}
The antenna was simulated using the following functions: \\
\begin{enumerate}
  \item singleDipole() $\rightarrow$ This function calculates the E- and H-field contributions from a single Hertzian dipole at a single point in space
  \item oneHertzianAllPoints() $\rightarrow$ This function calculates the E- and H-field contributions from a single Hertzian dipole to all points in space
  \item fullTraceAllPoints() $\rightarrow$ This function calculates the E- and H-field contributions from a whole meander line trace to all points in space
  \item MeanderTotalRadiation() $\rightarrow$ This function calculates the E- and H-field contributions of all Meander traces for all points in space
\end{enumerate} \\

The results were:

$$
E_{HPBW} = 117.30^{\circ}, ~~~~ H_{HPBW} = 151.68^{\circ}
$$

Which gives a directivity of:
$$
D = \frac{41253}{117.3^{\circ}\times 151.68^{\circ}} = 2.3184 = 3.65 \text{dBi}
$$


\section*{Conclusion}
There was a discrepancy between calculations and simulations of:
$$
4.4~\text{dBi} - 3.65~\text{dBi} = 0.75~\text{dBi}
$$

\noindent The radiation pattern can be seen by running \texttt{Meander\_Line\_Radiation\_Pattern\_3D.html} or by running the \texttt{Antenna\_sim.ipynb}

\end{document}